%
% reversi_160_costruzione_modello.tex
%
% This file is part of the reversi program
% http://github.com/rcrr/reversi-v2
%
% Author Roberto Corradini mailto:rob_corradini@yahoo.it
% Copyright 2026 Roberto Corradini. All rights reserved.
%
% License
%
% This program is free software; you can redistribute it and/or modify it
% under the terms of the GNU General Public License as published by the
% Free Software Foundation; either version 3, or (at your option) any
% later version.
%
% This program is distributed in the hope that it will be useful,
% but WITHOUT ANY WARRANTY; without even the implied warranty of
% MERCHANTABILITY or FITNESS FOR A PARTICULAR PURPOSE. See the
% GNU General Public License for more details.
%
% You should have received a copy of the GNU General Public License
% along with this program; if not, write to the Free Software
% Foundation, Inc., 59 Temple Place - Suite 330, Boston, MA  02111-1307, USA
% or visit the site <http://www.gnu.org/licenses/>.
%

%-----------------------------------------------------------------------------
% Capitolo 8: Costruzione del modello.
%
% Il dettaglio della terza fase della pipeline: i livelli DESIGN_MATRIX, ZED,
% GRADIENT, OPTIMIZE, GENERATE. L'ottimizzazione viene introdotta solo per
% nome (il solutore L-BFGS): il suo trattamento matematico completo è
% rimandato al capitolo 9, "L'algoritmo L-BFGS".
%-----------------------------------------------------------------------------

\chapter{Costruzione del modello}

Apre qui il dettaglio della terza fase della pipeline — \emph{Costruzione del
modello} — quella annunciata nella tabella dei 19 livelli del capitolo 5,
dove occupa le righe centrali della sequenza (i livelli 8–12:
\texttt{DESIGN\_MATRIX}, \texttt{ZED}, \texttt{GRADIENT}, \texttt{OPTIMIZE},
\texttt{GENERATE}). Le fasi precedenti
(Setup e Corpus e analisi delle feature) hanno preparato tutto ciò che questa
fase consuma: il feature set risolto, le posizioni di addestramento con i loro
valori di gioco, gli indici categorici di ogni feature e le mappe di peso
che associano a ogni configurazione di feature l'indice del suo peso nel
vettore $\bm{w}$.

Questa fase \emph{assembla} il problema di regressione e lo \emph{risolve}.
Nei cinque passi, in ordine:
\begin{itemize}
  \item \texttt{DESIGN\_MATRIX} costruisce la matrice di disegno $\bm{X}$: per
    ogni posizione, l'elenco degli indici dei pesi attivi;
  \item \texttt{ZED} trasforma i valori di gioco $\bm{y}$ nello spazio del
    predetto lineare, con la trasformazione logit (e le sue inverse);
  \item \texttt{GRADIENT} definisce la funzione di perdita $f(\bm{w})$ e il
    suo gradiente, confezionati nella closure \texttt{fg} che il solutore
    richiama a ogni iterazione;
  \item \texttt{OPTIMIZE} minimizza $f$ e produce il vettore di pesi
    ottimizzato $\bm{w}$;
  \item \texttt{GENERATE} assembla il vettore di pesi denso \texttt{w\_dense},
    che è il modello finale.
\end{itemize}
Il filo conduttore è semplice: i primi tre livelli descrivono \emph{il
problema} (dove stanno i pesi, verso quale target, e quale superficie di
errore scendere), il quarto lo \emph{risolve}, il quinto \emph{confeziona} la
soluzione nella forma che l'inferenza richiede. Ogni livello segue la divisione
di responsabilità già vista nel capitolo 5 — il worker
(\texttt{wrk\_*}, che gestisce cache, logging e pubblicazione sul contesto) e
l'helper (\texttt{hlp\_*}, che fa il calcolo di dominio) — e il suo contratto
di input/output è riassunto in tabella~\ref{tab:costruzione-contracti}:

\begin{table}[htbp]
\centering
\small
\begin{tabular}{lll}
\toprule
Livello & Consuma (dal contesto) & Pubblica (sul contesto) \\
\midrule
\texttt{DESIGN\_MATRIX} & indici principali, \texttt{iwmap} & \texttt{design\_matrix} ($M \times P$), il suo checksum \\
\texttt{ZED}           & valori di gioco $\bm{y}$, logit clipping & \texttt{z}, le closure \texttt{y2z}/\texttt{z2y}, il suo checksum \\
\texttt{GRADIENT}      & \texttt{design\_matrix}, \texttt{z}, $\bm{w}$, ridge & la closure \texttt{fg} \\
\texttt{OPTIMIZE}      & \texttt{fg}, $\bm{w}$ iniziale & $\bm{w}$ ottimizzato, \texttt{w\_checksum}, \texttt{opt\_info} \\
\texttt{GENERATE}      & $\bm{w}$, \texttt{wmap}, \texttt{iwmap} & \texttt{w\_dense}, \texttt{dense\_w\_checksum} \\
\bottomrule
\end{tabular}
\caption{I contratti dei cinque livelli della fase di costruzione del modello.}
\label{tab:costruzione-contracti}
\end{table}

Visto in blocco, ciò che questa fase assembla è un problema di
\emph{minimizzazione non lineare} di una funzione scalare
$f: \mathbb{R}^{N} \to \mathbb{R}$ — la perdita — il cui minimo fornisce il
vettore di pesi $\bm{w}^{*}$. Il problema è grande (centinaia di migliaia di
variabili) e strutturalmente sparso (ogni posizione del corpus tocca solo una
piccola frazione dei pesi), e queste due proprietà determinano le scelte dei
cinque livelli: rappresentare i dati in modo compatto, lavorare solo con la
funzione e il suo gradiente (mai con la hessiana), e sfruttare la sparsità a
ogni valutazione.

C'è infine un vincolo di \emph{qualità numerica} che lega
\texttt{GRADIENT} e \texttt{OPTIMIZE}: i metodi quasi-Newton, L-BFGS
compreso, presuppongono che la funzione e il gradiente forniti dal closure
\texttt{fg} siano \emph{consistenti} — cioè che il gradiente sia la
derivata esatta (a meno di arrotondamenti) della funzione che il solutore
sta minimizzando. Per questo le due quantità vengono calcolate nella stessa
valutazione, con la stessa precisione e sugli stessi indici: un gradiente
approssimato o incoerente farebbe deragare la ricerca lineare e
l'approssimazione della hessiana.

Una seconda trama attraversa i cinque livelli: la \emph{politica di cache},
che varia da livello a livello in funzione del costo di ricalcolo.
\texttt{DESIGN\_MATRIX} e \texttt{GENERATE} scaricano il proprio artefatto su
file binario e lo rilettono a una nuova run se il checksum di coerenza
concede: la matrice di disegno e il vettore denso sono gli artefatti più
costosi della fase. \texttt{ZED} e \texttt{GRADIENT} non usano la cache a
disco, perché il loro calcolo (aritmetica vettorizzata e generazione di una
closure) è trascurabile rispetto al resto. \texttt{OPTIMIZE} ha un
meccanismo a sé, i \emph{checkpoint} periodici, che non è una cache ma una
assicurazione: l'ottimizzazione è il passo più lungo della run e non può
permettersi di ricominciare da zero a ogni avvio.

Il flusso di dati della fase è lineare, e a ogni giunzione sta un checksum:
$\bm{X}$ e $\bm{z}$ entrano nel closure \texttt{fg}, \texttt{fg} entra nel
solutore, il vettore $\bm{w}$ ottimizzato esce dal solutore e entra
nell'espansione di \texttt{w\_dense}. La catena di checksum
(\texttt{design\_matrix\_checksum} $\to$ \texttt{zed\_checksum} $\to$
\texttt{w\_checksum} $\to$ \texttt{dense\_w\_checksum}) è ciò che rende la
fase insieme \emph{riproducibile} (stesso input, stesso output, verificabile
a posteriori) e \emph{riprendibile} (una run interrotta riparte dal primo
artefatto il cui checksum non torna).

\section{La matrice di disegno (\texttt{DESIGN\_MATRIX})}

Il livello \texttt{DESIGN\_MATRIX} costruisce la matrice di disegno $\bm{X}$,
l'ingrediente di input del problema di regressione. Essa non contiene valori
continui: è una matrice intera di \emph{indici}. Se $M$ è il numero di
posizioni nel corpus di addestramento e $P$ è il numero di istanze di feature
attive in ciascuna posizione, la forma di $\bm{X}$ è $(M, P)$ e ogni entry
$X_{ij}$ è l'indice, nel vettore di pesi $\bm{w}$, della configurazione di
feature attiva in $i$-esima posizione sulla $j$-esima colonna. Il predetto
lineare di una riga è allora un semplice \emph{gather}:

\begin{equation}
  \eta_i \;=\; \sum_{j=1}^{P} w_{\,X_{ij}} .
  \label{eq:design-predictor}
\end{equation}

Questa formulazione è la chiave di tutta l'efficienza del sistema: $\bm{X}$
riferisce solo \emph{dove} andare a prendere ogni contributo, e la somma
\eqref{eq:design-predictor} è l'unica operazione che li unisce. Non viene mai
materializzata una matrice float $M \times P$: ogni entry occupa 4 byte
(\texttt{uint32}) e il lavoro di memoria è tutto nel sommatorio indicizzato.

Il senso delle colonne si legge dal feature set risolto dal livello
\texttt{FEATURES}: ogni colonna di $\bm{X}$ è legata a un'istanza di feature
del set — la $j$-esima colonna raccoglie, per ogni posizione, l'indice del
peso dell'istanza $j$-esima (l'intercetto, le istanze dei pattern, le
istanze delle mobility, nell'ordine fissato dal set). Il numero di colonne
$P$ è quindi la somma degli \texttt{n\_instances} di tutte le feature, e
cambia da un feature set all'altro. Gli indici che la matrice contiene sono
gli indici \emph{principali}: quelli che, come visto nel capitolo della
funzione di valutazione, condensano in un solo representative le
configurazioni simmetriche, e sono proprio quelli per cui il vettore $\bm{w}$
ha un peso addestrato. Un ultimo dettaglio di tipo è rilevante: la mappa
\texttt{iwmap} è \texttt{int64} perché deve poter rappresentare la sentinella
$-1$ (configurazione non vista), mentre $\bm{X}$ è \texttt{uint32}, senza
sentinelle: la conversione è esplicita e avviene dopo la verifica che nessun
indice vi punti.

Merita un ricordo rapido di ciò che \texttt{iwmap} contiene, perché la
matrice di disegno è il suo consumatore diretto. Il livello \texttt{WMAPS}
partiziona l'intervallo $[0, N)$ degli indici di peso in una porzione per
feature (i confini sono \texttt{iwmap\_feature\_offset}) e, dentro ogni
porzione, assegna a ogni configurazione \emph{osservata} sopra la soglia di
frequenza un slot proprio; tutte le altre configurazioni di quella feature —
le rare, le mai viste — sono puntate a un unico slot di \emph{fallback}
della feature. È questa compressione a tenere $N$ sotto controllo: il
vettore di pesi addestrati è tanto lungo quanto le configurazioni che il
corpus ha effettivamente supportato, più un rappresentante per feature.

Per fissare le idee, un feature set con l'intercetto, un pattern a 4 istanze
(es. \texttt{EDGE}, i quattro bordi) e una mobility a 1 istanza produce
matrici con $P = 1 + 4 + 1 = 6$ colonne. Una riga — i numeri sono puramente
illustrativi — potrebbe essere:

\begin{center}
\small
\begin{tabular}{cccccc}
\toprule
intercetto & bordo 1 & bordo 2 & bordo 3 & bordo 4 & mobility \\
\midrule
0          & 7         & 1         & 12        & 3         & 42        \\
\bottomrule
\end{tabular}
\end{center}

dove ogni numero è l'indice, in $\bm{w}$, del peso da sommare per quella
colonna: il predetto lineare della posizione è
$w_0 + w_7 + w_1 + w_{12} + w_3 + w_{42}$. In un modello reale $P$ è
dell'ordine di una decina e $M$ di decine di migliaia, ma la struttura — e il
costo, $O(M \cdot P)$ — resta la stessa.

Va sottolineato il contrasto con la matrice di disegno \emph{classica} della
regressione lineare, che avrebbe $M$ righe e una colonna \emph{per ogni peso}
($N$ colonne, con $N$ dell'ordine di diverse centinaia di migliaia),
quasi tutta piena di zeri. Quella rappresentazione richiederebbe
$M \times N$ float, ovvero decine di gigaByte per un corpus di dimensioni
normali, e un prodotto matrice-vettore denso a ogni valutazione. La forma
adottata qui — una colonna per \emph{istanza di feature}, con entry che sono
\emph{indici} e non coefficienti — comprime la stessa informazione in
$M \times P$ interi a 4 byte, e sposta il costo dal prodotto denso al
sommatorio indicizzato \eqref{eq:design-predictor}, che tocca solo i pesi
attivi. È la stessa idea, in forma matriciale, per cui il gradiente che
vedremo al livello \texttt{GRADIENT} è non nullo solo sulle configurazioni
effettivamente presenti sulla scacchiera.

La costruzione riprende l'artefatto del livello \texttt{WMAPS}. Gli indici
principali delle feature (\texttt{pindexes}, prodotti dal livello
\texttt{INDEXES}) sono un array \texttt{uint32} di forma $(M, P)$ — per ogni
posizione e per ogni istanza di feature, l'indice della configurazione
attiva, ridotto alla sua classe di simmetria — e sono ancora
\emph{locali} a ogni feature; per convertirli
in indici globali nel vettore $\bm{w}$, a ogni riga si somma l'offset della
feature di appartenenza (\texttt{iwmap\_feature\_offset}) e si risolve
l'indice con la mappa inversa \texttt{iwmap} (configurazione $\to$ indice di
peso, con la sentinella $-1$ per le configurazioni non viste):

\begin{lstlisting}[language=python, caption={Il cuore di \texttt{hlp\_design\_matrix\_compute}
(estratto).}]
pindexes = ctx.hlp_indexes.indexes
n_instances_per_feature = [f.n_instances for f in ctx.feature_set.features]
col_offsets = np.repeat(ctx.iwmap_feature_offset[:-1], n_instances_per_feature)
pindexes_with_offsets = pindexes + col_offsets

X = np.empty(pindexes.shape, dtype=np.uint32)
np.take(ctx.iwmap.astype(np.uint32), pindexes_with_offsets, out=X)
\end{lstlisting}

Se anche un solo indice risolvesse sulla sentinella $-1$, il calcolo fallirebbe
immediatamente: è una garanzia di coerenza tra la mappa di peso e gli indici
del corpus. Il risultato è incapsulato nell'oggetto \texttt{LMDesignMatrix},
che accoglie la matrice $\bm{X}$ insieme al checksum delle mappe di peso che
l'hanno generata (\texttt{wmaps\_obj\_checksum}): quell'impronta è la chiave
di coerenza della cache, e la matrice stessa viene firmata con un suo
checksum SHA3-256 calcolato sui byte. La persistenza a disco usa il formato
binario comune (header con descrizione e versione, campi serializzati,
compressione LZ4 dell'intero file, file sidecar con il checksum): lo stesso
meccanismo che tutti i livelli a cache della pipeline condividono.

Il worker pubblica sul contesto la matrice nuda
(\texttt{ctx.design\_matrix}) e il suo checksum (SHA3-256), e sfrutta la
ricetta \emph{load-or-compute} della cache: una matrice già in cache viene
riletta da disco solo se la sua impronta di coerenza (il checksum delle mappe
di peso che l'hanno generata) coincide con quella della run corrente. Il
meccanismo ha una conseguenza pratica importante: se tra una run e l'altra
cambia qualunque cosa che modifichi le mappe di peso — il feature set, la
soglia di frequenza, il corpus — il checksum non coincide più e la matrice
viene silenziosamente ricalcolata, senza che l'utente debba cancellare nulla
a mano.

\section{Il target trasformato (\texttt{ZED})}

Il target della regressione non è il valore di gioco grezzo: il modello
logistico predice con la sigmoide, il cui rango è $(0, 1)$, mentre i valori di
gioco $y_i$ (il differenziale di dischi finale di ogni partita risolta) sono
interi nell'intervallo $[-64, 64]$. Prima di poterli regredire, vanno portati
nello spazio del predetto lineare — quello che il codice chiama lo spazio
\emph{affine} $z$ — con la trasformazione logit, nella sua forma
\emph{linearizzata} attorno al punto medio.

I valori \texttt{y} arrivano dal corpus come interi \texttt{int8} (il
differenziale di dischi finale di una partita su 64 caselle non può
eccederli); il vettore \texttt{z} che ne deriva è \texttt{float32}, lo stesso
tipo di precisione di tutto il modello. Con $c$ il parametro di \emph{logit
clipping} (dalla sezione \texttt{stat\_model} della configurazione), la
trasformazione e la sua inversa sono:

\begin{align}
  z_i &= \frac{1}{2} + y_i\;\frac{1 - 2c}{128} ,
  \label{eq:y2z}\\[4pt]
  y_i &= \frac{128}{1 - 2c}\left(z_i - \frac{1}{2}\right) .
  \label{eq:z2y}
\end{align}

La \eqref{eq:y2z} mappa l'intervallo intero $[-64, 64]$ suriettivamente
nell'intervallo $[c, 1-c] \subset (0, 1)$; la \eqref{eq:z2y} è la sua inversa
esatta, e verrà usata più avanti dalla fase di validazione per riportare le
predizioni del modello nello spazio dei valori di gioco. Il denominatore
$128$ non è arbitrario: è l'estensione completa del range dei valori di
gioco, da $-64$ a $+64$, e fa sì che la pendenza della mappa sia
$\frac{1-2c}{128}$ per unità di dischi — con $c$ piccolo, il target si
muove di meno di un centesimo per ogni disco di vantaggio, e il
punto di parità $y = 0$ cade esattamente al centro dell'intervallo, $1/2$.

In codice, le due
mappature sono due closure generate da una stessa fabbrica di funzioni,
parametrata su $c$:

\begin{lstlisting}[language=python, caption={Il cuore di \texttt{hlp\_zed\_fun\_factory}
(estratto).}]
a  = np.float32(0.5)
b  = np.float32((1.0 - 2.0 * logit_clipping) / 128.0)
def y_to_z(y):
    return a + np.float32(y) * b

a1 = np.float32(-64.0 / (1.0 - 2.0 * logit_clipping))
b1 = np.float32(128.0 / (1.0 - 2.0 * logit_clipping))
def z_to_y(z):
    return a1 + z * b1
\end{lstlisting}

Perché la forma affine? La trasformazione logit generale,
$z = \sigma^{-1}(p) = \log\bigl(p/(1-p)\bigr)$, non ha un'immagine a
intervallo finito: i valori estremi del gioco la farebbero esplodere, e la
sigmoide inversa introduce instabilità numerica proprio dove i residui
sarebbero più grandi. La forma affine \eqref{eq:y2z} è la sua versione
\emph{linearizzata}: conserva la geometria del problema (target dentro il
rango della sigmoide, monotono rispetto a $y$) con una mappa lineare,
invertibile in modo esatto e stabile a ogni passo.

Il parametro di clipping $c$ è il margine di sicurezza: i target
non si avvicinano mai ai bordi $0$ e $1$, dove la derivata della sigmoide —
e con essa il gradiente del prossimo livello — collasserebbe. Con $c$ piccolo
(ad esempio $0.01$), la mappa è pressoché lineare su tutto il suo dominio:
$y = +64$ va in $z = 0{.}99$, $y = 0$ resta esattamente in $z = 1/2$
(punto di perfetta simmetria: posizione equa), $y = -64$ va in $z = 0{.}01$.
Il guadagno numerico è concreto: tenendo i target lontani dai bordi, il
modello non è mai costretto a predire valori arbitrariamente vicini a $0$ o a
$1$ — cosa che richiederebbe un predetto lineare $|\eta|$ enorme e, con
esso, pesi di grandi dimensioni. Il training resta così in un regime in cui
la sigmoide è ben condizionata e il gradiente non si azzera né esplode.

La tabella riassume la corrispondenza per $c = 0.01$:

\begin{table}[htbp]
\centering
\small
\begin{tabular}{rrrrrr}
\toprule
$y$ & $+64$ & $+32$ & $0$ & $-32$ & $-64$ \\
\midrule
$z$ & $0.99$ & $0.745$ & $0.5$ & $0.255$ & $0.01$ \\
\bottomrule
\end{tabular}
\caption{Il target trasformato $z$ in funzione del valore di gioco $y$, per
\texttt{logit\_clipping} $c = 0.01$.}
\label{tab:zed-esempio}
\end{table}

Il vettore \texttt{z} di forma $(M,)$ e le due closure \texttt{y2z}/\texttt{z2y}
vengono pubblicati sul contesto (con il checksum di \texttt{z}): il calcolo è
vettorizzato e istantaneo, quindi qui non interviene la cache a disco, a
differenza del livello precedente. Il worker logga a scopo di verifica
l'intervallo effettivo del vettore prodotto (minimo e massimo di
\texttt{z}): deve coincidere con $[c, 1-c]$ a meno di arrotondamenti, e una
scostazione sarebbe il sintomo che il clipping configurato non è quello che
si credeva. Le closure non sono un dettaglio
implementativo: sono l'unico canale con cui il resto della pipeline traduce
tra i due spazi. Il livello di ottimizzazione lavora interamente nello spazio
$z$ (il target è $\bm{z}$, il predetto è la sigmoide, che vive in $(0,1)$);
la fase di validazione userà \texttt{z2y} per riportare le predizioni del
modello nel linguaggio dei valori di gioco, dove le metriche di errore hanno
senso per chi legge il report. Il checksum del vettore \texttt{z} non è solo
diagnosi: è uno dei campi con cui il livello \texttt{OPTIMIZE} verificherà
la coerenza del proprio checkpoint, così che una run ripresa da disco sia
certa di stare ottimizzando verso lo stesso target con lo stesso
\emph{design} di quando si era interrotta.

\section{La funzione di perdita e il suo gradiente (\texttt{GRADIENT})}

Con $\bm{X}$ e $\bm{z}$ a posto, il livello \texttt{GRADIENT} definisce il
problema da ottimizzare. Il modello logistico, per la $i$-esima posizione, è:

\begin{align}
  \eta_i &= \sum_{j=1}^{P} w_{\,X_{ij}} ,
  \label{eq:linear-predictor}\\
  \hat z_i &= \sigma(\eta_i) = \frac{1}{1 + e^{-\eta_i}} ,
  \label{eq:sigmoid-predict}
\end{align}
dove la sigmoide è calcolata in forma numericamente stabile
(\texttt{expit} di SciPy), per non sovraccaricare sugli $\eta$ estremi.

La funzione di perdita è l'errore quadratico medio sui residui tra predizione
e target, normalizzato sul numero di posizioni $M$, più il termine di
regolarizzazione di ridge (norma $L_2$ dei pesi), con coefficiente
$\lambda$ (il \texttt{ridge\_regularization} della configurazione):

\begin{equation}
  f(\bm{w}) \;=\;
  \frac{1}{2M}\sum_{i=1}^{M}\bigl(\hat z_i - z_i\bigr)^{2}
  \;+\; \frac{\lambda}{2M}\,\lVert\bm{w}\rVert^{2} .
  \label{eq:loss}
\end{equation}

La normalizzazione per $M$ fa sì che la scala di $f$ e dei suoi incrementi
sia indipendente dalla dimensione del corpus: gli stessi criteri di
convergenza valgono per un dataset piccolo e per uno grande. Il termine di
ridge tiene i pesi di dimensioni contenute (la motivazione è nel capitolo
della funzione di valutazione) e rende la superficie di errore più regolare.

Un punto va detto con franchezza: la composizione della sigmoide con
l'errore quadratico \emph{non} produce una funzione convessa, a differenza
della cross-entropy, e il minimo di \eqref{eq:loss} non è quindi garantito
convesso in modo formale. Ciò che rende il problema trattabile è che la
superficie è liscia, il suo gradiente è sempre disponibile e a buon mercato,
e la regolarizzazione di ridge tiene i pesi in un regime in cui la discesa
converge stabilmente a un ottimo locale di qualità — il ragionamento
completo, con lo sviluppo dell'hessiano, è nel capitolo della funzione di
valutazione. Per il solutore la conseguenza pratica è una sola: ciò che gli
serve, a ogni passo, sono esattamente il valore di \eqref{eq:loss} e il
gradiente che segue, e non nessuna informazione di curvatura esatta.

Il gradiente si ottiene con la regola della catena, applicata lungo la catena
$\bm{w} \to \eta_i \to \hat z_i \to f$. Poiché
$\frac{\partial \hat z_i}{\partial \eta_i} = \hat z_i(1-\hat z_i)$ e
$\frac{\partial \eta_i}{\partial w_k}$ conta quante volte l'indice $k$
compare nella riga $i$, si ha:

\begin{align}
  \frac{\partial f}{\partial w_k}
  &= \frac{1}{M}\sum_{i=1}^{M}
     \bigl(\hat z_i - z_i\bigr)\,\hat z_i\bigl(1-\hat z_i\bigr)\,
     \frac{\partial \eta_i}{\partial w_k}
  \;+\; \frac{\lambda}{M}\,w_k .
  \label{eq:grad-chain}
\end{align}

Il residuo della $i$-esima riga è $r_i = \hat z_i - z_i$; poiché la riga $i$
tocca i pesi solo in corrispondenza degli indici $X_{ij}$,
l'equazione \eqref{eq:grad-chain} diventa:

\begin{equation}
  \nabla f(\bm{w}) \;=\;
  \frac{1}{M}\sum_{i=1}^{M}
  r_i\,\hat z_i\bigl(1-\hat z_i\bigr)\,\mathbf{1}_i
  \;+\; \frac{\lambda}{M}\,\bm{w} ,
  \label{eq:gradient}
\end{equation}
dove $\mathbf{1}_i$ è il vettore indicatore della riga: ha $1$ nelle $P$
posizioni $X_{ij}$ e $0$ altrove. In forma compatta, con
$\bm{r} = \sigma(\bm{Xw}) - \bm{z}$ e $\bm{d} = \hat z \circ (1 - \hat z)$
(prodotto element-wise):

\begin{equation}
  \nabla f(\bm{w}) \;=\;
  \frac{1}{M}\,\bm{X}^{T}\bigl(\bm{d} \circ \bm{r}\bigr)
  \;+\; \frac{\lambda}{M}\,\bm{w} .
  \label{eq:gradient-compact}
\end{equation}

L'equazione \eqref{eq:gradient} è anche una dichiarazione di efficienza: per
ogni posizione, il gradiente è non nullo \emph{solo} sui pesi delle
configurazioni effettivamente attive. Il codice sfrutta questa sparsità
strutturale: anziché costruire la matrice $\bm{X}^{T}$, scatterizza i
contributi per riga direttamente negli indici, con \texttt{np.bincount}.
Ciascuna riga $i$ contribuisce al gradiente con lo stesso peso
$r_i\,\hat z_i(1-\hat z_i)$ in tutte le sue $P$ componenti, e il
\texttt{bincount} somma i contributi di tutte le righe che condividono lo
stesso indice di peso: il costo del gradiente è $O(M \cdot P)$, lo stesso del
predetto, e non dipende da $N$.

Il risultato del livello è la closure \texttt{fg} (\texttt{ctx.fg}): una
funzione che accetta il vettore di pesi corrente, come array \texttt{float32}
di $N$ componenti, e restituisce la coppia $(f, g)$ — il valore della perdita
come scalare \texttt{float32} e il gradiente come vettore \texttt{float32} di
$N$ componenti — pronta a essere richiamata a ogni iterazione del solutore. La closure cattura $\bm{X}$, $\bm{z}$ e il
coefficiente normalizzato $\alpha = \lambda/M$ una sola volta, e pre-alloca i
buffer di lavoro. La variante classica, vettorizzata in NumPy, è:

\begin{lstlisting}[language=python, caption={La closure \texttt{fg} della variante
classica (estratto di \texttt{hlp\_gradient\_compute}).}]
def fg(w):
    linear_predictor = np.sum(w[X], axis=1)
    zh  = sigmoid(linear_predictor)
    dzh = zh * (1.0 - zh)
    rn  = zh - z

    # Normalized loss (Mean Squared Error + Ridge regularization)
    f = 0.5 * (np.dot(rn, rn) / M + alpha * np.dot(w, w))

    # Normalized gradient
    g0 = np.bincount(X.ravel(),
                     weights=(dzh * rn)[:, None].ravel(),
                     minlength=N) / M
    g  = g0 + alpha * w
    return np.float32(f), g
\end{lstlisting}

Un dettaglio di ingegneria merita attenzione: la closure è richiamata dal
solutore a ogni iterazione, e ogni iterazione può valutarla più volte (la
ricerca lineare richiama la funzione nei punti di prova del passo). Perciò
ogni allocazione che non sia indispensabile va fatta \emph{una volta sola},
alla costruzione della closure: la variante classica pre-alloca fuori dal
closure il buffer di dimensioni $M \times P$ di float32 che il
\texttt{bincount} usa come pesi, e lo riutilizza a ogni chiamata
(azzerandolo via vista). All'interno del closure non resta che aritmetica
vettorizzata su buffer già in memoria.

Esiste una seconda variante, abilitata dalla configurazione
(\texttt{high\_performance\_fg}): un unico kernel Numba
(\texttt{\_hp\_fg\_kernel}) che percorre le righe di $\bm{X}$ in un solo
passaggio in streaming, partiziona le righe in $T$ blocchi contigui
(\texttt{hi\_perf\_blocks}, default $8$) e fa accumulare ogni blocco nel
proprio buffer privato di gradiente a 64 bit — senza condizioni di race
(Numba non ha atomic su CPU) e con risultato deterministico a prescindere dal
numero di thread, perché la partizione delle righe e l'ordine in cui i buffer
parziali vengono infine sommati non dipendono dalla parallelizzazione. Il
punto di progetto del kernel è il traffico di memoria: la variante classica
materializza a ogni valutazione l'array intermedio $M \times P$ del predetto
lineare (e la sua vista per il \texttt{bincount}), mentre il kernel in
streaming legge $\bm{X}$ una sola volta, in ordine di righe, e tiene tutto
lo stato di lavoro — il vettore $\bm{w}$, i buffer dei blocchi — residente
in cache L2. Le due varianti sono matematicamente identiche: differiscono
solo nell'ordine di arrotondamento a 32 bit (qualche ulp su $f$ e $g$), per
cui una run con l'una o con l'altra segue una traiettoria di ottimizzazione
lateralmente equivalente.

\section{L'ottimizzazione (\texttt{OPTIMIZE})}

Il livello \texttt{OPTIMIZE} risolve il problema

\begin{equation}
  \bm{w}^{*} \;=\; \operatorname*{arg\,min}_{\bm{w}} f(\bm{w}) ,
  \label{eq:optimize}
\end{equation}
con il solutore quasi-Newton \textbf{L-BFGS} (\emph{Limited-memory BFGS}),
implementato in \texttt{reversi.tools.opt\_lbfgs}. È la scelta naturale per
questo problema: il vettore di pesi conta centinaia di migliaia di
componenti, per cui un metodo di Newton pieno — che richiederebbe di
costruire e fattorizzare la matrice hessiana $N \times N$ — è inimmaginabile
in memoria; L-BFGS approssima l'inversa della hessiana con una piccola
storia di iterazioni e lavora solo con prodotti vettore, il che lo rende
adatto ai problemi grandi e strutturalmente sparsi come questo.

Il worker avvia il solutore passando la closure \texttt{fg} prodotta dal
livello precedente, il vettore di pesi iniziale — quello inizializzato a
zero dal livello \texttt{WMAPS}, salvo diverso riavvio da checkpoint — e i
comandi di convergenza
dalla sezione \texttt{optimization} della configurazione: il numero massimo di
iterazioni (\texttt{max\_iters}), la soglia di norma del gradiente
(\texttt{min\_grad}), i coefficienti delle condizioni di accettazione del
passo (\texttt{c1}, \texttt{c2}), il decremento minimo della funzione
(\texttt{min\_p\_fun\_decrease}) e la lunghezza $m$ della storia. Durante la
run il worker esegue checkpoint periodici su disco (\texttt{save\_every\_n}),
così che una interruzione — o una run successiva — possa riprendere
dall'ultima iterazione completata invece che da capo, e logging periodico
(\texttt{log\_every\_n}). Ciascun checkpoint salva, oltre ai pesi correnti,
l'intera memoria del solutore (la storia dei passi e dei gradienti che
approssimano la hessiana) e il contatore delle iterazioni completate, insieme
ai campi di coerenza che permettono alla run successiva di capire se quel
checkpoint è ancora valido: i checksum della matrice di disegno e del
target, il coefficiente di ridge e i parametri del solutore stessi. Una
run successiva può quindi riprendere \emph{esattamente} dal punto in cui la
precedente si era interrotta, e se il checkpoint risulta già convergente — o
ha esaurito \texttt{max\_iters} — l'ottimizzazione viene saltata
interamente e si riparte direttamente dai pesi finali; se la coerenza non
torna (la run è cambiata), il checkpoint viene scartato e si riparte da
capo. Al termine, il worker
pubblica sul contesto il vettore ottimizzato \texttt{ctx.w}, il suo checksum
e le informazioni di esito (numero di iterazioni, perdita finale, motivo di
arresto).

Il trattamento completo dell'algoritmo — come la storia delle iterazioni
approssima l'inversa della hessiana, la ricorsione a due cicli che calcola
la direzione di ricerca e la ricerca lineare che ne valida il passo — è
oggetto del prossimo capitolo, \emph{L'algoritmo L-BFGS}. Qui ci si limita a
stabilire che l'ottimizzazione \eqref{eq:optimize} è affidata a quel solutore
e che le quantità che gli occorrono a ogni passo sono esattamente le due
fornite dalla closure \texttt{fg}: il valore della perdita e il suo gradiente.

\section{La generazione del modello (\texttt{GENERATE})}

Il vettore $\bm{w}$ ottimizzato è il prodotto dell'apprendimento, ma non è
ancora il modello: contiene un slot per ogni configurazione
\emph{osservata} (sopra la soglia di frequenza) più i slot di fallback, e la
sua lunghezza $N$ non copre lo spazio completo delle configurazioni. Lo
schema sparso di $\bm{w}$ è un'esigenza del \emph{training}: le
configurazioni rare vengono delegate a un unico slot di fallback, e il
solutore deve poterle raggiungere attraverso la mappa \texttt{iwmap}.
L'\emph{inferenza} non ha questa macchina: l'engine che userà il modello deve
poter leggere il peso di una configurazione con un semplice indice diretto,
per ogni configurazione possibile, anche se rara o mai vista. Il livello
\texttt{GENERATE} risponde a questa esigenza: espande $\bm{w}$ nel vettore
\texttt{w\_dense} di lunghezza $K$ — il totale delle configurazioni
\emph{teoriche} di tutte le feature (l'ultimo offset di
\texttt{iwmap\_feature\_offset}, quindi $K \geq N$) — nel quale ogni
configurazione ha il suo peso, e che è la forma in cui il modello viene
distribuito e usato in inferenza. Lo scalamento è concreto: un pattern con
centinaia di migliaia di configurazioni teoriche ne avrà addestrate — con
sopra la soglia di frequenza — solo una frazione, e sarà l'espansione a
popolare il resto.

Il punto di partenza è la mappa \texttt{wmap} del livello \texttt{WMAPS},
che per ogni feature elenca le triple
\texttt{[feature\_id, config\_id, frequency]}: le righe con
\texttt{config\_id} $\geq 0$ identificano le configurazioni \emph{vere}
(osservate sopra la soglia di frequenza, con lo slot di peso proprio), le
altre sono configurazioni delegate al fallback. Da qui si ricava, per ogni
feature, l'elenco delle configurazioni con peso addestrato, e da lì la regola
di espansione, che dipende dal tipo di feature:

\begin{itemize}
  \item \textbf{Intercept}: la feature ha un solo peso; \texttt{w\_dense} lo
    replica diretto sulla sua intera porzione.
  \item \textbf{Mobility} (feature ordinale): le configurazioni osservate
    hanno pesi addestrati, le altre no; i pesi di \texttt{w\_dense} sono
    l'\emph{interpolazione lineare} dei pesi osservati sull'intero range delle
    configurazioni, con gli esterni valori clamped sul confine più vicino
    (l'interpolazione di \texttt{np.interp}, applicata su configurazioni
    garantite in ordine ascendente). Ad esempio, se del range di valori
    possibili ne sono stati osservati solo alcuni, un valore intermedio riceve
    la media lineare dei pesi dei due osservati più vicini, e un valore oltre
    l'osservato più estremo riceve il peso di quell'estremo: la scala
    ordinale della feature giustifica l'ipotesi che il peso vari
    \emph{regolarmente} da una configurazione all'altra.
  \item \textbf{Pattern} (feature categorica): le configurazioni
    \emph{simmetriche} ereditano il peso della loro configurazione
    \emph{principale} (il representative della classe di simmetria) — la
    stessa condivisione per cui, durante l'addestramento, tutte le
    configurazioni di una classe puntano a un unico slot di \texttt{w}, ora
    estesa a \emph{tutte} le configurazioni della classe; le configurazioni
    non viste e quelle di fallback, invece, ereditano la \emph{media pesata
    per frequenza} dei pesi delle configurazioni osservate di quel pattern.
    La media pesata è la stima plug-in del contributo atteso del pattern
    sotto la distribuzione di addestramento; il peso del fallback addestrato
    resta nel vettore $\bm{w}$ (dove tiene la perdita ben definita e
    ridimensiona le configurazioni rare), ma non è quello usato in
    inferenza, e viene solo registrato a log per trasparenza.
\end{itemize}

Il vettore denso è un array \texttt{float32} di $K$ componenti, e sarà
l'unico artefatto numerico del modello a uscire dalla pipeline: è quello che
il livello \texttt{SAVE} persisterà nel file di modello e che il livello
\texttt{CLIENT\_CHECK} verificherà essere di nuovo caricabile. La
conseguenza sul piano della coerenza è che la validità della cache di
\texttt{w\_dense} è agganciata al checksum di $\bm{w}$: qualsiasi
cambiamento che si ripercuota sull'ottimizzazione — anche solo un
divergere della traiettoria del solutore — invalida il vettore denso
cacheato e ne forza la rigenerazione, così che il modello distribuito sia
sempre la conseguenza esatta dei pesi ottimizzati della run corrente. Anche qui il
worker segue la ricetta \emph{load-or-compute}: \texttt{w\_dense} viene
scaricato su file binario (con compressione) e a una nuova run riletto se il
checksum del vettore di pesi ottimizzato da cui deriva coincide con quello
della run corrente. Sul contesto vengono pubblicati \texttt{ctx.w\_dense} e il
suo checksum.

\section{Consegna alla fase di validazione}

Un'ultima osservazione sul significato di ciò che questa fase produce. Il
vettore $\bm{w}$ ottimizzato è l'oggetto \emph{statistico}: la risposta,
sotto il vincolo della perdita \eqref{eq:loss}, alla domanda «quanto vale, in
media, una configurazione di feature?». Il vettore \texttt{w\_dense} è lo
stesso oggetto \emph{operativo}: la risposta resa consultabile in $O(1)$ per
ogni configurazione, incluse quelle che il corpus non ha mai fornito. La
qualità di entrambi si misurerà solo sulla fase successiva.

La fase è completa: il contesto porta \texttt{ctx.w} (i pesi ottimizzati) e
\texttt{ctx.w\_dense} (il modello denso), ciascuno con la propria impronta
di integrità. Il bilancio dei cinque livelli è il problema di regressione
integro e risolto: la matrice di disegno dice \emph{dove} stanno i pesi, il
target trasformato dice \emph{cosa} devono spiegare, la closure
\texttt{fg} dice \emph{come} misurare quanto ne sono lontani,
l'ottimizzazione ha trovato il punto in cui sono più vicini possibile, e la
generazione ha confezionato il risultato nella forma che l'engine consumerà.

Il livello successivo della pipeline, \texttt{VLD\_POSITIONS}, carica il set
di validazione e la fase di \emph{Validazione e analisi}
(\texttt{VALIDATE}, \texttt{ANALYTICS}) eseguirà il modello su posizioni mai
viste durante l'apprendimento: è là che si misurerà se i pesi prodotti da
questa fase generalizzano, e dove la closure \texttt{z2y} tornerà in scena
per riportare le predizioni dello spazio $z$ nello spazio dei valori di gioco.
